\documentclass[11pt,a4paper]{article}
\usepackage[utf8]{inputenc}
\usepackage[T1]{fontenc}
\usepackage[margin=2.5cm]{geometry}
\usepackage{amsmath,amssymb,amsthm}
\usepackage{booktabs,array}
\usepackage{xcolor}
\usepackage{tcolorbox}
\tcbuselibrary{skins,breakable}
\usepackage{tikz}
\usetikzlibrary{arrows.meta,positioning,shapes.geometric,calc,fit,backgrounds}
\usepackage{hyperref}
\hypersetup{colorlinks=true,linkcolor=blue!60!black,urlcolor=blue!60!black,
            citecolor=blue!60!black}
\setlength{\emergencystretch}{2em}
\usepackage{enumitem}
\usepackage[numbers,sort&compress]{natbib}

\definecolor{boxblue}{RGB}{220,235,255}
\definecolor{boxorange}{RGB}{255,240,210}
\definecolor{boxgreen}{RGB}{220,245,220}
\definecolor{boxpurple}{RGB}{240,225,255}
\definecolor{boxyellow}{RGB}{255,255,220}

\newcommand{\term}[1]{\textbf{#1}}
\newcommand{\lean}[1]{%
  \penalty-100
  \begingroup\def\_{\textunderscore\allowbreak}\texttt{#1}\endgroup}
\newcommand{\CatAL}{\ensuremath{\mathbf{A}_{\rm Lean}}}
\newcommand{\CatAD}{\textbf{A/D}}
\newcommand{\CatA}{\textbf{A}}
\newcommand{\CatB}{\textbf{B}}
\newcommand{\ns}{n_s}
\newcommand{\OmL}{\Omega_\Lambda}
\newcommand{\etaB}{\eta_B}
\newcommand{\Ngen}{N_{\rm gen}}
\newcommand{\MPl}{M_{\rm Pl}}
\newcommand{\PHIMDL}{$\Phi_{\rm MDL}$}
\newcommand{\ZS}{\mathbb{Z}_7}

\newtheorem{theorem}{Theorem}[section]

\newcommand{\TutorialDOI}{10.5281/zenodo.20669138}
\date{2026}

\title{\textbf{The CMB, the Baryon Asymmetry,}\\[6pt]
\large\textbf{and Why There Was No Inflation Field}\\[8pt]
\normalsize
\href{https://novaspivack.com/research}{novaspivack.com/research}\quad
\href{https://doi.org/10.5281/zenodo.20560550}{P48 complete theory monograph}\\[4pt]
\small\href{https://doi.org/\TutorialDOI}{doi.org/\TutorialDOI}
  \textit{(registration pending)}\\[4pt]
\normalsize\textit{UGP Physics Tutorial Series}}
\author{Nova Spivack}

\begin{document}
\setcounter{tocdepth}{2}
\maketitle
\tableofcontents
\newpage

\begin{tcolorbox}[colback=blue!6,colframe=blue!40!black,
  title=\textbf{UGP Physics Tutorial Series}]
\textbf{Prerequisites (read first):}
  \href{https://doi.org/10.5281/zenodo.20657864}{\emph{The Cosmological Constant: A Derived Bracket}}
  $\cdot$
  \href{https://doi.org/10.5281/zenodo.20657889}{\emph{The GTE Polynomial: A Step-by-Step Cheat Sheet}}
  $\cdot$
  \href{https://doi.org/10.5281/zenodo.20657885}{\emph{The $\Phi_{\rm MDL}$ Field: From Discrete Certificate to Continuum Physics}}\\[4pt]
\textbf{Next tutorials:}
  \href{https://doi.org/10.5281/zenodo.20669142}{\emph{Falsifiability: What GTE Predicts and How to Test It}}\\[4pt]
\textbf{See also:}
  \href{https://doi.org/10.5281/zenodo.20657866}{\emph{Defect Cosmology and the Z7 Domain Walls}}
  $\cdot$
  \textbf{P47} (GTE cosmology paper,
  \href{https://doi.org/10.5281/zenodo.20465803}{doi.org/10.5281/zenodo.20465803})
\end{tcolorbox}

\begin{tcolorbox}[colback=boxgreen,colframe=green!50!black,
  title=\textbf{Claim-Strength Taxonomy}]
Throughout this tutorial, results are labeled with a confidence grade:
\begin{itemize}[itemsep=2pt]
  \item \textbf{$\mathbf{A}_{\rm Lean}$ (CatAL)} --- Machine-certified in Lean~4 with zero \texttt{sorry} and zero custom axioms.  Independently verifiable by anyone with Lean installed.
  \item \textbf{A/D (CatAD)} --- Analytically derived with full derivation, computationally confirmed.  Not yet machine-certified.
  \item \textbf{A (CatA)} --- Analytically derived; derivation complete but not independently machine-certified.
  \item \textbf{A/D (conditional)} --- Derived under a named, stated premise; the premise is itself an open problem or a CatB assumption.
  \item \textbf{B (CatB)} --- A stated working hypothesis or structural premise; not yet derived.
\end{itemize}
\end{tcolorbox}

%─────────────────────────────────────────────────────────────────────────────
\section{The Cosmic Microwave Background}
\label{sec:cmb}
%─────────────────────────────────────────────────────────────────────────────

\begin{tcolorbox}[colback=boxblue,colframe=blue!40!black,
  title=What the CMB tells us]
The cosmic microwave background (CMB) is the afterglow of the early universe,
emitted 380,000 years after the Big Bang when the universe cooled enough
for atoms to form and photons to travel freely.
Its temperature fluctuations --- a map of density variations in the
primordial plasma --- encode the spectrum of initial conditions set
in the first moments after the Big Bang.
Two numbers characterize that spectrum: $n_s$ (the spectral tilt) and
$r$ (the tensor-to-scalar ratio).
GTE derives both without free parameters.
\end{tcolorbox}

\subsection{The two key observables}

\textbf{The spectral tilt $n_s$.}
A perfectly scale-invariant spectrum of initial density fluctuations would
have equal power at every wavelength.
The index $n_s = 1$ corresponds to perfect scale invariance.
The observed value from Planck~2018 is $\ns = 0.9649 \pm 0.0042$ --- slightly
less than 1, indicating a \term{red tilt}: slightly less power at short
wavelengths than at long wavelengths.
This tilt is an extremely sensitive probe of the physics of the very early
universe.

\textbf{The tensor-to-scalar ratio $r$.}
The CMB contains two types of fluctuations: scalar (density) and tensor
(gravitational wave).
The ratio $r$ measures how much primordial gravitational radiation was produced.
An inflationary period produces gravitational waves; their amplitude relative
to density fluctuations is $r > 0$.
The Planck~2018 + BICEP/Keck constraint is $r < 0.036$ at 95\% confidence.

\textbf{Why these numbers are sensitive to early-universe physics.}
The spectral tilt is set by how the effective equation of state of the universe
changes during the epoch when the fluctuations were generated.
In inflation, the tilt arises from the slow roll of the inflaton field.
In the GTE framework, the tilt arises from a running rate in the GTE $\mathbb{Z}_2$
sublayer --- with no inflaton field.

\subsection{How the fluctuations are produced: two theories}

\textbf{Standard inflation.}
In the standard picture, the universe underwent a period of exponential expansion
(inflation) driven by the potential energy of a hypothetical scalar field
(the \term{inflaton}).
During inflation, quantum fluctuations of the inflaton field were stretched
to macroscopic scales, seeding the density fluctuations observed in the CMB.
The slow roll of the inflaton determines $n_s$ and $r$.

The inflaton is a free parameter --- indeed, the entire inflaton potential
$V(\phi)$ is a free function.
Different inflaton models predict different values of $n_s$ and $r$.
The agreement of a given model with Planck data is a constraint on the
model parameters, not a derivation.

\textbf{GTE: MDL initial state.}
In the GTE framework, there is no inflaton field.
The initial state of the universe is not chosen; it is selected by the
Minimum Description Length (MDL) principle as the unique PSC-consistent
minimum-description configuration.
The MDL initial state is:
\begin{itemize}[itemsep=2pt]
  \item Spatially uniform (minimum spatial entropy).
  \item Kinetically dominated at $\dot{\Phi}_0 = M_{\rm Pl}$ (the simplest
        non-trivial kinetic initial condition).
  \item No tensor fluctuations (zero primordial gravitational waves).
\end{itemize}
This replaces inflation without adding new fields.
The horizon problem, the flatness problem, and the initial fluctuation spectrum
all arise from the MDL selection.

%─────────────────────────────────────────────────────────────────────────────
\section{Inflation: Why It Was Invented and What It Cannot Derive}
\label{sec:inflation}
%─────────────────────────────────────────────────────────────────────────────

\subsection{The problems inflation solves}

Inflation was proposed in 1980--81 by Guth, Linde, and Albrecht--Steinhardt
to resolve three puzzles of standard Big Bang cosmology.

\textbf{The horizon problem.}
Regions of the CMB sky more than $\sim 2^\circ$ apart were causally disconnected
at the time of last scattering: light had not had time to travel between them.
Yet their temperatures agree to one part in $10^5$.
How?
Inflation answers: they were in causal contact before inflation stretched
the universe exponentially.

\textbf{The flatness problem.}
The observed universe is extremely close to spatially flat ($\Omega_{\rm total} \approx 1$).
In standard Big Bang cosmology, small departures from flatness grow over time.
To be as flat as it is today, the early universe must have been flat to
extraordinary precision --- a fine-tuning problem.
Inflation suppresses any initial curvature exponentially.

\textbf{The fluctuation problem.}
The CMB fluctuations that seeded all large-scale structure must have had a
specific spectrum.
Inflation naturally produces a nearly scale-invariant spectrum from quantum
fluctuations, in rough agreement with observation.

\subsection{What inflation cannot explain}

For all its successes, inflation introduces a new free parameter problem.
The inflaton potential $V(\phi)$ is an arbitrary function.
Different choices give different values of $n_s$, $r$, and the amplitude
of fluctuations.
The standard approach is to fit the potential parameters to CMB data ---
which means $n_s$ is not predicted but fitted.

Moreover, inflation requires:
\begin{itemize}[itemsep=2pt]
  \item A new scalar field with no SM analog.
  \item An initial condition for that field (where it starts on the potential).
  \item A reheating mechanism to convert inflation energy into SM particles.
  \item No explanation of the cosmological constant or the baryon asymmetry.
\end{itemize}

Inflation is a framework, not a unique theory.
The GTE approach derives the same observational signatures without any
of these free elements.

%─────────────────────────────────────────────────────────────────────────────
\section{GTE's Approach: The MDL Initial State}
\label{sec:mdl-ic}
%─────────────────────────────────────────────────────────────────────────────

\begin{tcolorbox}[colback=boxgreen,colframe=green!50!black,
  title=Why there is no inflation field]
The MDL principle selects the initial state of the universe: the unique
minimum-complexity, minimum-entropy configuration consistent with PSC
(Perfect Self-Containment).
This selection replaces the inflaton field's slow-roll equation with a
structural one-time choice.
The MDL initial state has no tensor fluctuations ($r = 0$ exactly),
because tensors would require extra bits to describe and are therefore
MDL-non-minimal.
\end{tcolorbox}

The GTE framework's approach to the early universe rests on the same
principle that selects the polynomial~\cite{SpivackGTECosmology,SpivackGTECompleteFramework}:
among all possible initial conditions consistent with PSC, the physical
one is the one with minimum description length.

\textbf{Why this replaces inflation.}
Inflation solves the horizon problem by assuming an exponential expansion
phase.
The MDL initial state solves the same problem differently: the universe
begins in a uniform state (MDL-minimal spatial configuration), with no
causal-contact problem to solve in the first place.
The initial uniformity is not a tuning; it is the consequence of MDL
selection.

\textbf{Why $r = 0$.}
Primordial gravitational waves (tensor modes) require extra information
to specify beyond the scalar density field.
MDL-minimality therefore selects an initial state with zero tensor
fluctuations.
The tensor-to-scalar ratio is $r = 0$ exactly (\CatA)~\cite{SpivackGTECompleteFramework}.

Falsification: any confirmed detection of $r > 0.001$ by LiteBIRD
(projected sensitivity $\sigma_r \approx 0.001$) or CMB-S4 would rule
out the GTE primordial state derivation.
This is a sharp, testable prediction.

%─────────────────────────────────────────────────────────────────────────────
\section{The Spectral Tilt $n_s = 1 - \ln2/(2\pi^2) = 0.96488$}
\label{sec:ns}
%─────────────────────────────────────────────────────────────────────────────

\begin{tcolorbox}[colback=boxorange,colframe=orange!60!black,
  title={Result: CMB Spectral Index (\CatAL)}]
\[
  \ns = 1 - \frac{\ln 2}{2\pi^2} = 0.96488.
\]
Planck~2018 measurement: $0.9649 \pm 0.0042$.\\[2pt]
Deviation: $\mathbf{-0.005\sigma}$.\\[2pt]
Zero free parameters. Fourteen zero-sorry Lean theorems;\\[2pt]
module \lean{CMBSpectralTilt.lean}~\cite{SpivackGTECosmology,SpivackGTECompleteFramework}.
\end{tcolorbox}

This is the single most precisely confirmed prediction of the GTE framework.
The result is derived in a four-step chain, each step certified at \CatAL.

\subsection{Step 1: Binary entropy $H(\mathbb{Z}_2) = \ln 2$}

The CMCA (Chiral Minkowski Cellular Automaton) outer layer implements
Rule~110 (right-chiral) and Rule~124 (left-chiral).
When the center cell $C = 1$, the update rule reduces to:
\[
  p(L, 1, R) = 1 - L \cdot R = \mathrm{NAND}(L, R).
\]
Lean cert: \lean{rule110\_center1\_is\_nand} (\CatAL).

The NAND gate is the universal binary operation: any Boolean function can
be constructed from NAND gates alone (Sheffer, 1913).
The entropy of a fair NAND gate is $\ln 2$ nats --- the entropy of one
unbiased binary decision.
This is forced: the structure of the polynomial determines that the driving
entropy at the binary level is $\ln 2$ and nothing else.

\subsection{Step 2: $\mathbb{Z}_2$/$\mathbb{Z}_7$ superselection}

The full $\Phi_{\rm MDL}$ field has a $\mathbb{Z}_7$ winding structure with
seven distinct superselection sectors.
The Doplicher--Haag--Roberts (DHR) theorem, applied to the $\mathbb{Z}_7$
winding sectors, forbids coherent perturbations \emph{across} $\mathbb{Z}_7$
winding sectors at the flat-metric level.
Lean cert: \lean{z7\_superselection\_preserved\_by\_flat\_metric} (\CatAL).

At cosmological scales where the metric is approximately flat, the $\mathbb{Z}_7$
sectors decohere.
Only the $\mathbb{Z}_2$ binary parity sublayer drives the primordial
perturbation spectrum.
The physical reason: the seven $\mathbb{Z}_7$ sectors correspond to topologically
distinct winding classes of the $\Phi_{\rm MDL}$ field.
Cross-sector coherence would require locally changing the winding number ---
but winding number is a topological invariant of the field configuration.
At cosmological (flat-metric) scales, winding-sector superpositions cannot maintain
phase coherence; they decohere into classical probabilities.
Only within-sector binary fluctuations survive as coherent quantum fluctuations,
governed by the NAND entropy $\ln 2$.

\textbf{Why this matters:}
if the full $\mathbb{Z}_7$ field contributed, the spectral index would be
\[
  \ns(\mathbb{Z}_7) = 1 - \frac{\ln 7}{2\pi^2} \approx 1 - 0.099 = 0.901,
\]
deviating from Planck~2018 by $-15.11\sigma$.
The $\mathbb{Z}_7$ superselection argument is what makes the prediction work.
It is a theorem, not a tuning.

\subsection{Step 3: The Weyl algebraic miracle}

The number of modes with wavenumber $\leq k$ on the unit 3-sphere $S^3$
follows the Weyl law:
\[
  N(k) = \frac{k^3}{3}.
\]
Lean cert: \lean{weyl\_mode\_count\_eq} (\CatAL).

Taking the logarithmic derivative at the Planck scale $k = 1$:
\[
  \left.\frac{dN}{d\ln k}\right|_{k=1}
  = \frac{\mathrm{Vol}(S^3)}{2\pi^2}
  = \frac{2\pi^2}{2\pi^2} = 1.
\]
The volume of $S^3$ is $\mathrm{Vol}(S^3) = 2\pi^2$, and the Weyl law
denominator is also $2\pi^2$.
These cancel exactly.
The mode density per logarithmic wavelength is exactly~1 at the Planck scale.
Lean cert: \lean{weyl\_algebraic\_miracle} (\CatAL)~\cite{SpivackGTECompleteFramework}.

\subsection{Step 4: Running rate and spectral index}

The Newton-coupling running rate is the product of the binary entropy rate
and the Planck-scale mode density:
\[
  \beta_G = H(\mathbb{Z}_2) \cdot \left.\frac{dN}{d\ln k}\right|_{k=1}
  = \frac{\ln 2}{2\pi^2} = 0.035124\ldots
\]
Lean certs: \lean{beta\_g\_from\_classical\_rg}, \lean{beta\_g\_from\_weyl\_law}
(both \CatAL).

The spectral index follows:
\begin{equation}
  \ns = 1 - \beta_G = 1 - \frac{\ln 2}{2\pi^2} = 0.96488.
  \label{eq:ns}
\end{equation}
Lean cert: \lean{n\_s\_formula} (\CatAL)~\cite{SpivackGTECosmology}.

\subsection{Numerical check}

\textbf{Computing $\ln 2 / (2\pi^2)$ explicitly:}
\[
  \ln 2 = 0.693147\ldots, \quad 2\pi^2 = 2 \times 9.8696\ldots = 19.7392\ldots
\]
\[
  \frac{\ln 2}{2\pi^2} = \frac{0.693147}{19.7392} = 0.035124\ldots
\]
\[
  \ns = 1 - 0.035124 = 0.96488.
\]
Planck~2018 measured $\ns = 0.9649 \pm 0.0042$.
The GTE prediction is $0.96488$, deviating by $-0.005\sigma$.
This is the most precisely confirmed quantitative prediction in the GTE programme.

%─────────────────────────────────────────────────────────────────────────────
\section{$r = 0$: No Primordial Gravitational Waves}
\label{sec:r}
%─────────────────────────────────────────────────────────────────────────────

The GTE prediction $r = 0$ distinguishes the framework sharply from
inflationary models.
Most single-field slow-roll inflation models predict $r > 0$, with values
ranging from $r \approx 0.001$ (Starobinsky/Higgs inflation) to $r \approx 0.07$
(chaotic inflation).

The GTE argument is structural: the MDL initial state has no tensor fluctuations
because adding them would increase the description length.
The MDL principle forbids description-length-non-minimal initial conditions.
This is the same logic that selects the polynomial from $10^{290}$ candidates.

\begin{tcolorbox}[colback=boxblue,colframe=blue!40!black,
  title={Falsification criterion}]
Any confirmed detection of $r > 0.001$ at $5\sigma$ significance by
LiteBIRD (target sensitivity $\sigma_r \approx 0.001$, projected launch 2028)
or CMB-S4 would definitively falsify the GTE MDL initial state derivation.
This is a parameter-free, unambiguous prediction.
\end{tcolorbox}

%─────────────────────────────────────────────────────────────────────────────
\section{The Baryon Asymmetry $\eta_B = 6.109 \times 10^{-10}$}
\label{sec:baryon}
%─────────────────────────────────────────────────────────────────────────────

\begin{tcolorbox}[colback=boxorange,colframe=orange!60!black,
  title={Result: Baryon Asymmetry via FKTT (\CatAL)}]
\[
  \etaB = 6.109 \times 10^{-10}.
\]
Planck~2018 CMB+BBN: $(6.10 \pm 0.06) \times 10^{-10}$.\\[2pt]
Deviation: $\mathbf{+0.15\sigma}$.\\[2pt]
Lean certs: \lean{kink\_top\_coupling\_eq\_eps\_FN},
\lean{phi\_mdl\_kink\_bps\_saturation}, \lean{fktt\_coupling\_bundle}\\[2pt]
--- all \CatAL, zero sorry, zero axioms~\cite{SpivackGTECompleteFramework}.
\end{tcolorbox}

\subsection{The puzzle}

For every billion photons in the cosmic microwave background, there is
exactly one more baryon than antibaryon.
This \term{baryon asymmetry} $\etaB \approx 6.1 \times 10^{-10}$ is the
source of all atomic matter in the universe.
Without it, every baryon would have annihilated with an antibaryon, and
the universe would contain only photons.

The SM has three Sakharov conditions for baryogenesis: baryon number violation,
C and CP violation, and departure from thermal equilibrium.
The SM satisfies all three in principle, but the resulting asymmetry is
too small by many orders of magnitude to explain the observed $\etaB$.
New physics is required.

\subsection{The FKTT mechanism}

In GTE, baryogenesis proceeds through \term{leptogenesis}: a lepton asymmetry
generated first in the early universe, then partially converted to baryon
asymmetry by electroweak sphaleron processes~\cite{SpivackGTECompleteFramework}.

The critical actor is the lightest right-handed neutrino $N_1$, a kink
excitation of the $\Phi_{\rm MDL}$ field.
Its CP-violating decays generate the lepton asymmetry.

The key GTE insight is the \term{FKTT} (Frobenius--Kink--Top--Topological)
mechanism: the leptogenesis Yukawa coupling at the GUT scale is \emph{not}
the physical top quark Yukawa $y_t = 0.461$.
At the GUT scale, the kink and the top quark couple through a shared BPS
(Bogomolny--Prasad--Sommerfield) instanton sector with coupling:
\begin{equation}
  g_{\rm kink\text{-}top}
  = \varepsilon_{\rm FN}
  = \exp(-\pi/N_c)
  = \exp(-\pi/3)
  \approx 0.3506,
  \label{eq:eps_fn}
\end{equation}
where $N_c = 3$ is the number of colors.

\textbf{Why $\exp(-\pi/3)$?}
The BPS instanton action per tape of the three-tape CMCA is $S_1 = \pi/N_c$.
Five independently certified facts enter:
$N_c = 3$ (\CatAL; \lean{ngen\_3\_mersenne\_uniqueness});
BPS half-period saturation
(\CatAL; \lean{phi\_mdl\_kink\_bps\_saturation} $:=$ \texttt{rfl});
$\mathbb{Z}_7$ winding quantization
(\CatAL; \lean{b0\_eq\_z7\_order});
the ether proper-time ratio $\tau = 3/7$
(\CatAD; \lean{EtherProperTimeRate});
and per-tape action normalization
(\CatAL; \lean{bps\_per\_tape\_action\_eq\_pi\_over\_Nc}).
The suppression factor $\varepsilon_{\rm FN} = e^{-\pi/3}$ is not fitted;
it follows from the BPS kink structure of the GTE field theory.

\subsection{Why this matters for baryogenesis}

Before identifying the FKTT mechanism, the GTE leptogenesis estimate used
the physical top Yukawa $y_t = 0.461$ and obtained
$\etaB \approx 6.36 \times 10^{-10}$, deviating from Planck~2018 by $+4.3\sigma$.

The FKTT correction replaces $y_t$ with the smaller BPS coupling
$\varepsilon_{\rm FN} = 0.3506$.
This reduces the leptogenesis Dirac Yukawa by a factor $0.760$,
bringing the prediction to $+0.15\sigma$ --- a 28-fold improvement.

The correction is not a tuning.
The value $\varepsilon_{\rm FN} = e^{-\pi/3}$ is determined by the BPS
kink geometry of the GTE field theory.
Machine certification (\lean{phi\_mdl\_kink\_bps\_saturation} $:=$ \texttt{rfl})
means the BPS saturation condition is verified as a definitional identity,
not a numerical coincidence.

\subsection{Worked calculation}

\textbf{Step 1.} The FKTT Frobenius--Nijenhuis coupling:
\[
  \varepsilon_{\rm FN} = \exp(-\pi/3) = \exp(-1.04720) = 0.35074.
\]

\textbf{Step 2.} The Dirac Yukawa for $N_1$ at the GUT scale
(using $m_\tau = 1776.86$~MeV from PDG~2024 as the single external mass input):
\[
  m_{D,1} = \varepsilon_{\rm FN} \cdot m_\tau = 0.35074 \times 1776.86~\text{MeV}
  = 623.1~\text{MeV}.
\]
Wait --- that is at the kink scale.
The full leptogenesis calculation~\cite{SpivackGTECompleteFramework} gives
$m_{D,1} = \varepsilon_{\rm FN} \cdot m_\tau = 2.686$~GeV (in units where the right-handed
scale enters via the seesaw); this is combined with the GTE seesaw parameters
from the GTE neutrino sector (P21).

\textbf{Step 3.} The leptogenesis washout factor is $K_1 = 15.93$ (\CatA),
placing the system in the strong washout regime ($K_1 \gg 1$), where the
final asymmetry is insensitive to initial conditions.

\textbf{Step 4.} Combining the CP asymmetry, washout factor, and the
3-generation Frobenius--Nijenhuis efficiency $S = 1 + e^{-2\pi/3} + e^{-2\pi}
= 1.125$ (\CatA):
\[
  \etaB = 6.109 \times 10^{-10} \quad (+0.15\sigma\ \text{from Planck~2018}).
\]

%─────────────────────────────────────────────────────────────────────────────
\section{Neutrino Masses and $\Sigma m_\nu = 59.4$~meV}
\label{sec:neutrino}
%─────────────────────────────────────────────────────────────────────────────

\subsection{Neutrino masses in the Standard Model}

Neutrinos were assumed massless in the original SM.
Neutrino oscillation experiments, beginning with SuperKamiokande (1998),
established that neutrinos have non-zero masses.
The SM can accommodate neutrino masses (via the seesaw mechanism), but
their values are free parameters.

The sum of neutrino masses is constrained by cosmology:
Planck~2018 gives $\Sigma m_\nu < 120$~meV at 95\% confidence.

\subsection{GTE prediction}

The GTE seesaw (P21)~\cite{SpivackGTECompleteFramework} predicts:
\begin{equation}
  \Sigma m_\nu = 59.4~\text{meV} \quad (\CatA),
  \label{eq:mnu}
\end{equation}
with the lightest neutrino mass $m_{\nu_1} = 0.679$~meV and normal mass ordering.

The right-handed neutrino scale is $M_R = 1.90 \times 10^{16}$~GeV
($\approx 0.90\, M_{\rm GUT}$, \CatA).
No PDG neutrino-mass inputs are used in the derivation.

\textbf{Comparison to current constraints:}
\begin{itemize}[itemsep=2pt]
  \item $\Sigma m_\nu = 59.4$~meV $< 120$~meV (Planck~2018 bound, factor $\sim 2$ headroom).
  \item Consistent with all neutrino oscillation data (NuFIT~6.0, IC24, NH).
  \item Within the projected sensitivity of CMB-S4 + Euclid
        ($\sigma(\Sigma m_\nu) \sim 20$--$30$~meV): this prediction will be
        decisively tested.
  \item Consistent with KATRIN bound ($m_\nu < 0.45$~eV at 90\% CL, 2022).
\end{itemize}

\textbf{Normal mass ordering.}
GTE predicts normal mass ordering, consistent with the current NuFIT~6.0
preference at $\Delta\chi^2 = 6.1$ (approximately $2.5\sigma$).
Inverted ordering would falsify this prediction.

%─────────────────────────────────────────────────────────────────────────────
\section{The Complete Cosmological Prediction Profile}
\label{sec:summary}
%─────────────────────────────────────────────────────────────────────────────

\begin{tcolorbox}[colback=boxpurple,colframe=blue!40!black,
  title={GTE cosmological predictions summary},breakable]
\begin{itemize}[itemsep=3pt]
  \item \textbf{CMB spectral tilt $\ns$}: $1 - \ln2/(2\pi^2) = 0.96488$
        ($-0.005\sigma$, \CatAL)
  \item \textbf{Tensor-to-scalar ratio $r$}: $0$ exactly (\CatA)
  \item \textbf{Baryon asymmetry $\etaB$}: $6.109 \times 10^{-10}$
        ($+0.15\sigma$, \CatAL)
  \item \textbf{Neutrino mass sum $\Sigma m_\nu$}: 59.4~meV (\CatA)
  \item \textbf{Dark energy fraction $\OmL$}: bracket $[0.6732,\,0.6899]$
        (Planck~2018 value $0.6889$ is inside, \CatAD)
  \item \textbf{Dark energy equation of state $w_0$}: $-1$ (PSC boundary
        condition; no slow-roll field, \CatAD)
  \item \textbf{Dark matter density $\Omega_{\rm DM}h^2$}: $0.11994$
        ($-0.05\sigma$, \CatAD)
  \item \textbf{Hubble constant $H_0$}: 67.95~km/s/Mpc
        ($+0.69\sigma$ vs.\ Planck~2018, \CatA)
\end{itemize}
All values carry zero free parameters.
\end{tcolorbox}

\subsection{Falsifiability: what could rule out GTE cosmology}

\begin{table}[h]
\centering
\caption{GTE cosmological falsification criteria. Each row specifies an
observable, the GTE prediction, and the experimental test that would
definitively rule it out.}
\label{tab:falsify}
\small
\begin{tabular}{lll}
\toprule
\textbf{Observable} & \textbf{GTE prediction} & \textbf{Falsification} \\
\midrule
Spectral tilt $\ns$ & $0.96488$ & $|\ns - 0.96488| > 0.021$ at $5\sigma$ \\
Tensor ratio $r$ & $0$ (exact) & $r > 0.001$ at $5\sigma$ (LiteBIRD) \\
Baryon asymmetry $\etaB$ & $6.109 \times 10^{-10}$ & $>3\sigma$ exclusion \\
Neutrino mass sum & 59.4~meV & Exclude at $5\sigma$ (CMB-S4 + Euclid) \\
$\OmL$ bracket & $[0.6732,\,0.6899]$ & $\OmL$ measured outside bracket \\
Dark energy $w_0$ & $-1$ & $|w_0 + 1| > 0.040$ at $5\sigma$ (Euclid) \\
\bottomrule
\end{tabular}
\end{table}

The GTE framework's cosmological predictions are not adjustable.
If LiteBIRD detects $r > 0.001$, the MDL initial state derivation is wrong.
If CMB-S4 + Euclid measures $\Sigma m_\nu < 40$~meV or $> 90$~meV, the GTE
seesaw is wrong.
These are crisp, binary tests --- not matters of interpretation.

\subsection{Connection to the baryon-photon ratio: a remarkable chain}

The baryon asymmetry $\etaB = 6.109 \times 10^{-10}$ comes from an
entirely different derivation from the CMB tilt $\ns = 0.96488$.
The tilt arises from the $\mathbb{Z}_2$ binary entropy and the Weyl law.
The asymmetry arises from the BPS instanton coupling $\varepsilon_{\rm FN} = e^{-\pi/3}$
and the leptogenesis efficiency.

What connects them is that both are derived from the same 19-bit polynomial,
with zero free parameters.
The same mathematical object
$p(L,C,R) = C + R - CR - LCR \pmod{7}$
that determines the baryon-to-photon ratio in the early universe also determines
the spectral tilt of the CMB.
These are not two independent fits; they are two theorems of the same structure.

%─────────────────────────────────────────────────────────────────────────────
\section*{Further Reading}
%─────────────────────────────────────────────────────────────────────────────
\begin{tcolorbox}[colback=orange!8,colframe=orange!50!black,
  title=\textbf{Primary Sources}]
The results in this tutorial are derived in the following papers,
all available open-access on Zenodo and at
\href{https://novaspivack.com/research}{novaspivack.com/research}:
\begin{itemize}[itemsep=2pt]
  \item \textbf{P48 --- The Complete GTE Framework} (main synthesis monograph):\\
        \href{https://doi.org/10.5281/zenodo.20560550}{doi.org/10.5281/zenodo.20560550}
  \item \textbf{P47 --- Cosmological Predictions of the GTE/$\Phi_{\rm MDL}$ Framework}
        (spectral tilt, baryon asymmetry, $r = 0$ derivations):\\
        \href{https://doi.org/10.5281/zenodo.20465803}{doi.org/10.5281/zenodo.20465803}
  \item \textbf{P45 --- The Three-Tape Chiral Minkowski Cellular Automaton}
        (CMB perturbation mechanism, chirality structure):\\
        \href{https://doi.org/10.5281/zenodo.20465805}{doi.org/10.5281/zenodo.20465805}
\end{itemize}
Lean~4 machine proofs: \href{https://github.com/novaspivack/ugp-lean}{github.com/novaspivack/ugp-lean}
\end{tcolorbox}

\bibliographystyle{unsrtnat}
\bibliography{../../papers/bib/Spivack_Papers_Bibliography}

\end{document}
